\section{就业冲击分层：哪些岗位先变}

白领通缩不是均匀发生的。Coding/Agent 会先影响那些任务数字化程度高、输出可验证、流程可拆解的岗位。它不一定马上消灭整个职业，但会先重组岗位内部的任务比例。

\subsection{岗位冲击矩阵}

\noindent\begin{longtable}{p{0.22\textwidth}p{0.34\textwidth}p{0.35\textwidth}}
\toprule
岗位/任务 & 先被压缩的部分 & 更难替代的部分 \\
\midrule
初级程序员 & 样板代码、简单 bug、脚本、接口封装 & 架构判断、复杂系统理解、上线责任。 \\
数据分析 & 清洗、SQL、报表、可视化初稿 & 指标定义、业务解释、因果判断。 \\
产品/运营 & 文档、竞品总结、流程配置 & 用户洞察、取舍、组织协调。 \\
研究助理 & 文献整理、实验脚本、结果汇总 & 问题提出、实验设计、失败解释。 \\
管理协调 & 状态同步、会议纪要、排期 & 冲突解决、资源取舍、责任承担。 \\
\bottomrule
\end{longtable}

\begin{importantbox}{职业变化的真实形态}
AI 不一定一次性替代岗位，而是先替代任务。岗位会被重组：低判断、高重复、可验证的部分被自动化；高责任、高上下文、高取舍的部分变得更重要。
\end{importantbox}

\subsection{本章小结}

就业冲击应按任务颗粒度分析。最先被压缩的是数字化、重复性、可验证任务；最难替代的是责任、判断、组织和价值取舍。

